\section{Facet variations and limiting source laws}\label{sec:variations}

Fix a prescribed canonical full maximum, or a canonical joint maximum
$(U,\alpha)$ with $\alpha<L$. Use the selections of
\cref{thm:selection-prescribed}. Write $U_i$ for a selected polygon,
$h_i$ for its support, $A_i=f_{U_i}$, and $q_i$ for its sampled
complete barrier. Its projection and middle interval are $I_i,J_i$.
All supports are bounded by $B$ in absolute value, the mesh tends
to zero, and the penalty coefficient $\eta_i$ tends to zero.
We take each polygonal variation first, with $i$ fixed.

\subsection{Moving a floating side and an endpoint}

For a nonaxis, nontop sampled normal $s$, let $\ell_{i,s}$ be the
outer facet length, $\ell^{\rm w}_{i,s}$ its part over
$I_i\setminus J_i$, and $\ell^{\rm m}_{i,s}$ its part over $J_i$.
Let $\tau_{i,s}$ be the length of the positive barrier graph over
$J_i$ carried by the inner line with this normal. The terminal first
line is counted when $\alpha<L$; an unused normal has $\tau_{i,s}=0$.
Distinct source lines have distinct slopes, so their intersection
points do not create ambiguity in these lengths.

\begin{lemma}[Finite facet variation]\label{lem:variation-finite}
With $w_{i,s}$ the total persistent weight at $s$, put
$b_{i,s}=4B\eta_iw_{i,s}$ and $\beta_i=\sum_s b_{i,s}\le4B\eta_i$.
Then
\begin{equation}\label{eq:variation-facet}
 \ell^{\rm w}_{i,s}\le\tau_{i,s}+b_{i,s}.
\end{equation}
At a positive vertical end face the outward derivatives of the
unpenalized score are $e_{i,R}-C_{i,R}$ and $e_{i,L}-C_{i,L}$, where
\begin{equation}\label{eq:variation-pressures}
 \begin{gathered}
 Q_{i,\pm}=A_i(j_{i,\pm})+q_i(j_{i,\pm}),\\
 C_{i,R}=(3Q_{i,+}-Q_{i,-})/4,\qquad
 C_{i,L}=(3Q_{i,-}-Q_{i,+})/4,\\
 e_{i,R}=A_i(r_i),\qquad e_{i,L}=A_i(l_i).
 \end{gathered}
\end{equation}
For all end faces, including zero ones,
\begin{equation}\label{eq:variation-endpoint-upper}
 e_{i,R}\le(C_{i,R}+b_{i,0})_+,\qquad
 e_{i,L}\le(C_{i,L}+b_{i,\pi})_+.
\end{equation}
\end{lemma}

\begin{proof}
Move the supporting line at $s$ outwards by $\varepsilon>0$, keeping
all other grid halfplanes. An old attaining point remains for each
unchanged support, so those supports stay fixed. For a positive
facet the moved support rises by $\varepsilon$ for sufficiently
small $\varepsilon$. The projection and $J_i$ are unchanged.
The gained outer strip has area
$\varepsilon\ell^{\rm w}_{i,s}+O(\varepsilon^2)$. The gained
forbidden area has area
$\varepsilon\tau_{i,s}+O(\varepsilon^2)$: an exposed inner
segment moves by its normal displacement, and finitely many line
intersections, floor contacts, and window endpoints contribute
quadratic errors. Thus the score derivative is
$\ell^{\rm w}_{i,s}-\tau_{i,s}$. The penalty derivative has absolute
value at most $b_{i,s}$. Maximality gives
\eqref{eq:variation-facet}; at zero facet length it is automatic.

Move a positive right end face from $r_i$ to $r_i+\varepsilon$.
Nonaxis supports and $q_i(x)$ stay fixed, whereas the middle
endpoints move by $\varepsilon/4,3\varepsilon/4$. Differentiating
the three integrals gives
\[
 e_{i,R}+\tfrac14(A_i+q_i)(j_{i,-})
          -\tfrac34(A_i+q_i)(j_{i,+})=e_{i,R}-C_{i,R}.
\]
Moving the left face outwards moves the middle endpoints by
$-3\varepsilon/4,-\varepsilon/4$ and gives $e_{i,L}-C_{i,L}$.
If an end face has height zero, its retained last sloping halfplane
and the floor already impose the endpoint. Relaxing the axis
constraint then leaves the body unchanged. This proves precisely
the positive-part bounds \eqref{eq:variation-endpoint-upper}.
\end{proof}

\subsection{The four neighboring lines}

\begin{lemma}[Neighbor exposure]\label{lem:variation-neighbors}
At an interior first source $t$ with source spacing $\delta$, put
$f=h_i(t)$, $g=h_i(t+L)$, $c=\cos\delta$, $s=\sin\delta$,
and $T=\tan(\delta/2)$. Let
\[
 \ell^{\rm circ}=
 \frac{h_i(t-\delta)+h_i(t+\delta)-2cf}{s}\ge0.
\]
The full exposed first-wall length satisfies
\begin{equation}\label{eq:variation-neighbor}
 \tau\le\tan\delta(|q^+|+T)+(2T-\ell^{\rm circ})_+,
 \qquad q^+=f+\frac{h_i(t+L+\delta)-cg}{s}-1.
\end{equation}
The corresponding second-source bound has the roles of the two
normal families interchanged.
\end{lemma}

\begin{proof}
Parameterize the first inner ray by
$z(v)=(f-1)u(t)+v u^\perp(t)$, $v\le v_0=g-1$.
To avoid both neighboring open quadrants, an exposed point must lie in
\[
 [\min(v_-,v_+),v_0]\ \cup\ [v_{\rm lo},v_{\rm hi}],
\]
where solving the four neighboring line inequalities gives
\[
 \begin{aligned}
 v_-&=\frac{h_i(t+L-\delta)-1-(f-1)s}{c},&
 v_+&=\frac{h_i(t+L+\delta)-1+(f-1)s}{c},\\
 v_{\rm lo}&=\frac{h_i(t+\delta)-1-(f-1)c}{s},&
 v_{\rm hi}&=\frac{(f-1)c-h_i(t-\delta)+1}{s}.
 \end{aligned}
\]
Indeed either a neighboring companion alternative holds, putting
the point in the first interval, or both neighboring first alternatives
must hold, putting it in the second. At zero the neighboring quadrant
is below the floor; at a proper terminal endpoint it is dominated
by the whole terminal first wall. These neighbors remain valid.
Subtracting gives $v_{\rm hi}-v_{\rm lo}=2T-\ell^{\rm circ}$.
If $d_-=(gc-h_i(t+L-\delta))/s$ and
$d_+=(h_i(t+L+\delta)-gc)/s$, then $d_-\le d_+$ and
\[
 v_0-v_-=\tan\delta(f+d_--1+T),\qquad
 v_0-v_+=\tan\delta(-q^++T).
\]
The first interval therefore has length at most
$\tan\delta(|q^+|+T)$. Floor clipping only shortens the exposed
pieces, proving the bound. The second-source calculation is the
same calculation in the interchanged orthonormal coordinates.
\end{proof}

In a full grid $\ell^{\rm circ}=\ell_{i,t}$; with additional
outer normals $\ell_{i,t}\le\ell^{\rm circ}$. Uniform support
and support-slope bounds imply, for a fixed constant $K$,
\begin{equation}\label{eq:variation-coarse}
 \tau_{i,t},\tau_{i,t+L}\le K\delta,
 \qquad
 \ell_{i,s}\le K\delta+b_{i,s}+\ell^{\rm m}_{i,s}
\end{equation}
at nonterminal source normals. Normals used only for circumscription
satisfy $\ell^{\rm w}_{i,s}\le b_{i,s}$.
The terminal first source instead satisfies
$\tau_{i,\alpha}\le |J_i|/\sin\alpha$; this mass is retained.
There is no terminal companion source, since its two-wall minimum
is dominated by the terminal first wall.

\begin{lemma}[Vanishing middle errors]\label{lem:variation-middle}
On every compact normal arc separated from $\theta_c$ and $L$,
$\sum_s\ell^{\rm m}_{i,s}\to0$.
\end{lemma}

\begin{proof}
Concavity and the bound $0\le A_i\le1$ give uniform roof-slope
bounds near $J_U$, because this interval stays a positive distance
from the projection endpoints. The roofs converge uniformly there.
On $[j_-+\varepsilon,j_+-\varepsilon]$, their one-sided slopes
converge uniformly to the middle slope $m$: bound each derivative
between its two secant slopes to points a fixed distance away and
then pass to the uniform limit. Thus every middle facet there has
normal arbitrarily close to $\theta_c$. The two remaining boundary
strips have total arclength at most a fixed constant times
$\varepsilon+o(1)$. Let $i\to\infty$ and then
$\varepsilon\downarrow0$.
\end{proof}

\begin{lemma}[The tilted middle junctions]\label{lem:variation-middle-junctions}
If a canonical full or joint maximum has $m\ne0$, then
$A'_-(j_-)>m>A'_+(j_+)$. Its middle facet has projection exactly $J_U$.
\end{lemma}

\begin{proof}
At its normal $n_c=u(\theta_c)=(-m,1)/\sqrt{1+m^2}$ and $x\in J_U$,
\[
 \begin{gathered}
 h_U(\theta_c)-1-\ip{(x,0)}{n_c}
 =A(x)/\sqrt{1+m^2}-1\le-\delta_m,\\
 \delta_m=1-(1+m^2)^{-1/2}>0.
 \end{gathered}
\]
There is therefore an open normal arc $\Omega$ around $n_c$ on
which this deficit is at most $-\delta_m/2$, uniformly for $x\in J_U$.
Its walls exclude every positive forbidden point above $J_U$.
Changing only supports in $\Omega$ by less than $\delta_m/4$
leaves the complete middle barrier unchanged; this also covers
the terminal wall if its normal belongs to $\Omega$.

If, for example, $A'_+(j_+)=m$, concavity makes all upper normal
cones at wing points $x\downarrow j_+$ tend to $\{n_c\}$.
Choose such a point $p=(x,A(x))$ with $A(x)<1$ and its normal cone
inside $\Omega$, and adjoin $p+\varepsilon e_y$ with its downward
segment. Taking the convex hull raises the exterior roof near $x$.
Outside $\Omega$, the positive support gap of $p$ has a positive
minimum, so a sufficiently small $\varepsilon$ leaves those
supports unchanged. Inside $\Omega$ the barrier remains unchanged.
The score strictly increases, a contradiction. A point below height
one exists: for $m<0$ the roof decreases to the right, and for
$m>0$ a matching right slope forces $A(j_+)<1$.
The left junction uses the same argument, with the signs reversed.
Concavity supplies the nonstrict inequalities initially, so both
must be strict. A facet extending outside $J_U$ would violate one
of them.
\end{proof}

\subsection{Actual curvature densities and corner velocities}

The upper surface-area measure of a convex cap is
$\sigma=h''+h$ in the distributional sense; for a polygon it is
$\sum_s\ell_s\delta_s$. Uniform support convergence gives weak
convergence of these measures, since for a smooth test function
$\varphi$ their integrals are $\int h(\varphi''+\varphi)$.
Their total masses are bounded by the perimeter of the fixed box.

\begin{theorem}[Regular curvature and shadow laws]\label{thm:variation-curvature}
For every prescribed canonical full maximum the upper curvature
measure has a bounded density away from $\theta_c,L$.
For a prescribed proper joint maximum, remove also $\alpha$;
on unused arcs away from $\theta_c$ the density is zero.
There is no singular continuous curvature, and the only possible
interior atoms are at these removed normals.

Put $F(t)=h_U(t)$, $G(t)=h_U(t+L)$ and, on their regular intervals,
\begin{equation}\label{eq:variation-velocities}
 p=F'-G+1,\qquad q=G'+F-1,\qquad
 u=F''+F,\qquad v=G''+G.
\end{equation}
Then $F,G\in W^{2,\infty}_{\rm loc}$ there and
\begin{equation}\label{eq:variation-controls}
 \begin{gathered}
 p'=u-1-q,\qquad q'=v-1+p,\qquad p\le1,\quad q\ge-1,\\
 0\le u\le\kappa(q),\qquad0\le v\le\kappa(p),\qquad
 \kappa(z)=\max\{|z|,(1+|z|)/2\}.
 \end{gathered}
\end{equation}
At almost every angle where both supports are regular,
\begin{equation}\label{eq:variation-shadow}
 \begin{array}{ll}
 p>0, q>0\ \Longrightarrow\ u=0, v\le1/2,\\
 p<0, q<0\ \Longrightarrow\ v=0, u\le1/2.
 \end{array}
\end{equation}
If the inner corner has negative height, $u=v=0$; at height zero,
$u,v<1$ almost everywhere.
\end{theorem}

\begin{proof}
Sum \eqref{eq:variation-coarse} on a compact regular normal interval.
The penalty sum tends to zero and the middle sum tends to zero by
\cref{lem:variation-middle}. The result bounds its limiting curvature
by $K$ times its angular length. Testing nonnegative continuous
functions proves measure domination, hence a bounded density.
The remaining finite set of normals can carry only atoms.
Saturation supplies the assertion on unused arcs.

For the sharper bound, combine
\eqref{eq:variation-neighbor} with
$\ell_{i,t}\le\tau_{i,t}+b_{i,t}+\ell^{\rm m}_{i,t}$.
If $\ell_{i,t}\ge2T$, the positive part vanishes; otherwise move
its $-\ell_{i,t}$ term to the left. Since $\tan\delta=\delta+O(\delta^3)$
and $2T=\delta+O(\delta^3)$, the two cases give
\[
 \ell_{i,t}\le\delta\kappa(q^+_{i,t})+O(\delta^2)
                 +b_{i,t}+\ell^{\rm m}_{i,t}.
\]
The constant is uniform in the boxed polygons. Support functions
are uniformly semiconvex: $h_i''\ge-h_i\ge-B$.
Uniform convergence therefore makes nearby one-sided support
secants converge to the limiting derivative at each differentiability
point. Their common bound gives convergence in local $L^1$.
Since $\kappa$ is Lipschitz, summing the last inequality on an
interval gives $\sigma_F\le\kappa(q)\,\dd t$. Lebesgue
differentiation yields $u\le\kappa(q)$ almost everywhere.
The interchanged source families give $v\le\kappa(p)$.
The bounded-density equation $h''=\rho-h$ gives the stated
$W^{2,\infty}$ regularity. Differentiating
\eqref{eq:variation-velocities} gives the two differential equations.
The support points $Fu+F'u^\perp$ and
$Gu^\perp-G'u$ lie in $U$; comparison with the companion
supporting halfplanes gives $F'\le G$ and $-G'\le F$,
hence $p\le1$ and $q\ge-1$.

Here is the additional finite calculation for shadowing. Put
$d^+=(h_i(t+\delta)-cf)/s$, $d^-=(cf-h_i(t-\delta))/s$,
and use the corresponding secants $d_\pm$ at $t+L$.
Let $p_\pm=d^\pm-g+1$ and $q_\pm=f+d_\pm-1$.
In a full grid these are its one-sided velocities; additional
outer normals may shorten the circumscribed facets.
If $p_+>T,q_+>T$, its whole
first inner ray lies inside the next quadrant: the two thresholds
on that ray exceed its endpoint $g-1$ by
$p_+-T$ and $\tan\delta(q_+-T)$.
Its exposed first-source length is zero. If also
$q_-+T\ge\tan\delta(p_-+T)$, the previous and next quadrants
cover the companion ray outside
$[-d_--T,T-d_+]$. Its exposed length is at most
$(2T-(d_+-d_-))_+\le(2T-\ell_{i,t+L})_+$.
On a compact subset of $\{p>0,q>0\}$ all these inequalities hold
eventually. The facet variation then gives
\[
 \ell_{i,t}\le b_{i,t}+\ell^{\rm m}_{i,t},\qquad
 \ell_{i,t+L}\le T+b_{i,t+L}+\ell^{\rm m}_{i,t+L}.
\]
Sum and pass to the limit as before, using $T/\delta\to1/2$.
Interchanging the local families and reversing the neighboring
triple gives the negative-negative assertion. For proper angles
this argument takes place inside the visited open source interval;
on unused regular arcs the densities vanish already. It does not
reflect the partial objective.

Finally set $c=(F-1)u(t)+(G-1)u^\perp(t)$, so
$c'=pu+qu^\perp$. On a compact negative-height corner interval
all sampled corners are eventually negative. Their quadrants have
no positive exposed source; the facet and middle error bounds
therefore force zero limiting curvature.
On the level set $c_y=0$, absolute continuity gives
$c_y'=c_y''=0$ almost everywhere. Put $z=c_x'$, $s=\sin t$,
$k=\cos t$. Then
\[
 p=zk,\quad q=-zs,\quad su+kv=s+k-2z.
\]
The sharp bounds and $\kappa(a)\le|a|+1/2$ imply
$su+kv\le|z|+(s+k)/2$. This excludes $z\le0$.
Nonnegativity gives $z\le(s+k)/2$, hence
$|p|,|q|\le(1+\sqrt2)/4$ and
$u,v\le(5+\sqrt2)/8<1$.
\end{proof}

\subsection{Endpoint complementarity and source balance}

Define $Q_\pm=A(j_\pm)+N_\alpha(U)(j_\pm)$ and
$C_R=(3Q_+-Q_-)/4$, $C_L=(3Q_--Q_+)/4$.
For a full maximum take $\alpha=L$.

\begin{lemma}[Endpoint complementarity]\label{lem:variation-complementarity}
For every prescribed maximum,
$e_R=(C_R)_+$ and $e_L=(C_L)_+$.
\end{lemma}

\begin{proof}
Trim the right end to $r_i-\varepsilon$. The retained roof is unchanged;
all supports decrease, so the new barrier is at most the old barrier.
Keeping the old barrier as an upper comparison on the new window gives
\[
 P_{k_i,\alpha}(U_i\cap\{x\le r_i-\varepsilon\})-P_{k_i,\alpha}(U_i)
 \ge\varepsilon(C_{i,R}-e_{i,R})+o(\varepsilon).
\]
Concavity bounds the trimming distance by
$\varepsilon\sqrt{1+(w_i-\varepsilon)^{-2}}$:
the possible vertical change is at most the horizontal change divided
by $w_i-\varepsilon$, using the secant from the opposite endpoint.
Eventually $w_i\ge1$, so this is at most $\sqrt5\varepsilon$ for
$\varepsilon\le1/2$. The penalty change is at most
$4\sqrt5B\eta_i\varepsilon+O(\eta_i\varepsilon^2)$.
Maximality gives $e_{i,R}\ge C_{i,R}-4\sqrt5B\eta_i$.
The direct left trim gives the same left inequality. Together with
\eqref{eq:variation-endpoint-upper} these give
$e_{i,Q}=(C_{i,Q})_++o(1)$ for $Q=R,L$.

End heights require a separate convergence argument. Roof slopes
over $J_i$ are between $-4$ and $4$ eventually. Facets in a
sufficiently short punctured arc of either axis consequently have
zero middle length. Equation \eqref{eq:variation-coarse} bounds
their total length by $K(a+\delta_i)+\beta_i$ on an arc of
length $a$. Weak convergence of the full surface-area measures
therefore gives
\[
 \limsup e_{i,R}\le e_R\le\liminf e_{i,R}+Ka.
\]
The first inequality tests the closed axis singleton; the second
tests its open $a$-neighborhood. Let $a\downarrow0$.
The left argument is identical. Barriers and interior roof values
converge uniformly, and the middle endpoints converge, so
$C_{i,Q}\to C_Q$. Passing to the limit proves the lemma.
\end{proof}

\begin{lemma}[Positive pressures]\label{lem:variation-endpoint-pressure}
For every canonical full or joint maximum,
\[
 A(j_-)>1/2,\quad A(j_+)>1/2,\qquad
 e_R=C_R>0,\quad e_L=C_L>0.
\]
\end{lemma}

\begin{proof}
The higher middle endpoint has height one by
\cref{lem:selection-top}. The exact one-angle estimate proved in
\cref{lem:variation-low-height} gives $P_\alpha(U)<4/5$ if the
other endpoint is at most $1/2$, contradicting
$P_\alpha(U)\ge M/2>4/5$.
If $Q_+\ge1$ and $C_R\le0$, then $Q_-\ge3Q_+$;
complementarity gives $e_L=C_L\ge2Q_+\ge2$, impossible.
If the remaining pressure is nonpositive, then
$Q_+\ge3Q_-$ and $1\ge e_R=C_R\ge2Q_-$, forcing
$A(j_-)\le1/2$, again impossible. Exchanging sides proves the
opposite tilt. Both pressures are positive, so complementarity
removes their positive parts.
\end{proof}

\begin{theorem}[The limiting source measures]\label{lem:variation-sources}
Let $\omega$ be the actual normal-length measure of the upper boundary
over the exterior quarters, omitting the vertical end faces and the
horizontal top. The finite source measures
$\nu_i=\sum_s\tau_{i,s}\delta_s$ have subsequences converging weakly
to $\omega$. For full maxima this measure has bounded densities
on both open source quarters. For proper maxima it has bounded
densities on $(0,\alpha)$ and $(L,L+\alpha)$ and at most one atom,
at $\alpha$. It has no atom at the companion endpoint $L+\alpha$,
the axes, or the top, and no mass on unused open arcs.
\end{theorem}

\begin{proof}
The mesh bounds and the retained terminal bound give weakly compact
finite source measures. Equation \eqref{eq:variation-facet} gives
$\omega\le\nu$ for any joint weak limit of source and outer
wing measures. To identify the outer limit, test an angular function
$\varphi$ supported away from the axes and top. Its graph integral is
\[
 \int_{I_i\setminus J_i}
 \varphi\bigl(\arg(-A_i'(x),1)\bigr)\sqrt{1+A_i'(x)^2}\,\dd x.
\]
Concave roofs converge uniformly on interior intervals and their
derivatives converge almost everywhere. The angular cutoff bounds
the integrand even near steep slopes; dominated convergence and
shrinking the endpoint strips give the actual wing measure.
The coarse source bound excludes extra floating atoms at omitted
normals. The terminal first normal is included, so its facet mass
is retained in this identification.

For the reverse inequality use right horizontal erosion
$E^R_\varepsilon=U_i\cap(U_i-\varepsilon e_x)$.
Its roof is $\min(A_i(x),A_i(x+\varepsilon))$ and
\[
 h_{E^R_\varepsilon}(s)\le h_i(s)-\varepsilon(\cos s)_+.
\]
The virtual supports on the right dominate the actual eroded
barrier. Only first-source walls move in this comparison, including
the whole terminal first wall. Their finite barrier-area derivative
is $-\sum_{0<s<L}\tau_{i,s}\cos s$.
The exterior roof derivative is
$-\sum_{0<s<L}\ell^{\rm w}_{i,s}\cos s$.
Include the end sliver and the moving-window terms to obtain
\[
 \Delta P_i\ge\varepsilon(D_{i,R}-e_{i,R}+C_{i,R})+o(\varepsilon),
 \qquad D_{i,R}=\sum_{0<s<L}(\tau_{i,s}-\ell^{\rm w}_{i,s})\cos s.
\]
For distance control, put $\lambda=\varepsilon/w_i$. For $p\in U_i$,
both $(1-\lambda)p+\lambda(l_i,0)$ and this point plus
$\varepsilon e_x$ lie in $U_i$ by convexity. Thus erosion moves
points by at most $(\diam U_i/w_i)\varepsilon$, uniformly bounded
by a fixed multiple of $\varepsilon$.
Penalized maximality now gives
\[
 -\beta_i\le D_{i,R}\le e_{i,R}-C_{i,R}+O(\eta_i).
\]
Direct left erosion gives the same bound with weight $-\cos s$ on
the second quarter. Positive endpoint pressure and complementarity
make both limiting defect moments zero. Their weights are strictly
positive on their open quarters; at the first terminal atom the
weight is $\cos\alpha>0$. The source mesh estimates exclude atoms
at the zero-weight top endpoint. The nonnegative defect $\nu-\omega$
therefore vanishes on both quarters, proving $\nu=\omega$.
\end{proof}

\subsection{The stationary length and projection budgets}

Let $T_{\rm w}$ be the length of the height-one top face lying
outside $J_U$, and set $L_{\rm w}=\omega((0,\pi))+T_{\rm w}$.
Vertical end faces are excluded from this length.

\begin{corollary}[Stationary wing identity]\label{cor:variation-green}
Every canonical full or joint maximum satisfies
\begin{equation}\label{eq:variation-green}
 2P_\alpha(U)=L_{\rm w},\qquad
 \max_{x\in J_U}N_\alpha(U)(x)\le(e_R+e_L)/2.
\end{equation}
\end{corollary}

\begin{proof}
Put $O=\int_{I\setminus J}A$, $N=\int_JN_\alpha(U)$ and denote
barrier endpoint values by $n_\pm$. On an outer graph segment,
$h\,\dd s=(A-xA')\,\dd x$. Integration by parts on the two wings gives
\[
 2O=\int h\,\dd\omega+T_{\rm w}+re_R-le_L+j_-A_- -j_+A_+.
\]
On each finite positive inner graph segment the support distance is
$h_i-1$, so the same calculation gives
\[
 2\int_{J_i}q_i=\int(h_i-1)\,\dd\nu_i
                    -j_{i,-}q_i(j_{i,-})+j_{i,+}q_i(j_{i,+}).
\]
Uniform convergence and weak source convergence give this identity
in the limit. Subtract and use $\nu=\omega$:
\[
 2P_\alpha(U)=L_{\rm w}+re_R-le_L+j_-Q_- -j_+Q_+.
\]
The last two terms equal $lC_L-rC_R$; pressure equality cancels
all boundary terms.

Put $S=A_-+A_+$, $E=e_R+e_L$ and $Z=n_-+n_+$.
Top localization makes the wings monotone towards their middle
endpoints, so $\int|\cos s|\,\dd\omega=S-E$.
For each finite barrier,
$\operatorname{TV}_{J_i}(q_i)=\int|\cos s|\,\dd\nu_i$.
Lower semicontinuity of total variation gives
$\operatorname{TV}_{J}(N_\alpha(U))\le S-E$.
Continuity gives $2\max_JN_\alpha(U)-Z\le\operatorname{TV}_{J}(N_\alpha(U))$,
while the endpoint equations give $E=(S+Z)/2$.
Combining these proves the height bound.
\end{proof}

\begin{corollary}[Positive-barrier projection]\label{cor:variation-projection}
Every canonical full or joint maximum satisfies
\begin{equation}\label{eq:variation-projection}
 |\{x\in J_U:N_\alpha(U)(x)>0\}|\le w_U/2-T_{\rm w}.
\end{equation}
\end{corollary}

\begin{proof}
The horizontal projection of a finite source segment of length
$\tau_{i,s}$ is $\tau_{i,s}\sin s$. Their interiors are disjoint
parts of the barrier graph, so
\[
 |\{x\in J_i:q_i(x)>0\}|=\int\sin s\,\dd\nu_i.
\]
The outer projection budget similarly gives
$\int\sin s\,\dd\omega_i=w_i/2-T_{i,\rm w}$.
The limiting measure $\omega$ therefore has sine moment
$w_U/2-T_{\rm w}$, also obtained directly from its actual outer graph.
Uniform barrier convergence gives lower semicontinuity of the
positive-set length: each compact subset on which the limiting
barrier is positive is eventually contained in the finite positive
set. Apply inner regularity and $\nu_i\rightharpoonup\omega$.
\end{proof}

These identities use weak limits of finite source lengths.
They do not require those lengths to converge to the ordinary
arclength of the limiting niche: arbitrarily shallow positive
pieces may persist in the source measure after their heights vanish.
